\documentclass{ximera}
\title{Hyperbolic functions}
\begin{document}
\begin{abstract}
The hyperbolic functions, their definitions in terms of exponentials, their basic identities and derivatives, and the inverse hyperbolic functions with their domains and derivatives. Textbook §3.11.
\end{abstract}
\maketitle

Let's talk about hyperpbolic functions next. Recall that we can get the sin and cos functions from a circle. If instead we looked at a hyperbola, we would get the hyperbolic functions.
\begin{image}
\includegraphics[width=0.5\textwidth]{b48a5324-93dc-45ee-b7c2-dc0b5706bda6-101_481_608_1766_758}
\end{image}
\begin{image}
\includegraphics[width=0.5\textwidth]{b48a5324-93dc-45ee-b7c2-dc0b5706bda6-102_503_685_425_718}
\end{image}

They're hard to remember, so here are their definitions:
\[
\begin{array}{c|c||c|c}
\sinh (x) & \frac{e^{x}-e^{-x}}{2} & \csch(x) & \frac{1}{\sinh (x)} \\
\cosh (x) & \frac{e^{x}+e^{-x}}{2} & \sech(x) & \frac{1}{\cosh (x)} \\
\tanh (x) & \frac{\sinh (x)}{\cosh (x)} & \coth(x) & \frac{1}{\tanh (x)}
\end{array}
\]

Here is what $\sinh (x)$ and $\cosh (x)$ look like
\begin{image}
\includegraphics[width=0.45\textwidth]{b48a5324-93dc-45ee-b7c2-dc0b5706bda6-102_387_524_1330_535}
\end{image}
\begin{image}
\includegraphics[width=0.45\textwidth]{b48a5324-93dc-45ee-b7c2-dc0b5706bda6-102_385_524_1330_1061}
\end{image}

If we're lucky, we'll see some applications of these functions later in the term, but for now let's look at some of the properties of these functions.
\[
% CHECKED-FIX: was \sinh (x+y)=\sinh (x) \cosh (x)+\cosh (x) \sinh (y)  (identity typo, not an \answer)
\begin{array}{ll}
\sinh (-x)=-\sinh (x) & \cosh (-x)=\cosh (x) \\
\sinh (x+y)=\sinh (x) \cosh (y)+\cosh (x) \sinh (y) & \cosh ^{2}(x)-\sinh ^{2}(x)=1 \\
\cosh (x+y)=\cosh (x) \cosh (y)+\sinh (x) \sinh (y) & 1-\tanh ^{2}(x)=\sech^{2}(x)
\end{array}
\]

We can also compute the derivatives of the hyperbolic functions fairly easily using their original definitions using exponential functions.

\begin{example}\label{example:9-2}
Let $f(x)=\cosh (x)$, what is its derivative?
% work: cosh(x)=(e^x+e^{-x})/2, so f'(x)=(e^x-e^{-x})/2=sinh(x)
\[
f^{\prime}(x)=\frac{d}{dx}\left(\frac{e^{x}+e^{-x}}{2}\right)=\answer{\sinh(x)}
\]
\end{example}

Here are all the derivatives for the hyperbolic functions.

% work: tanh' = (cosh^2-sinh^2)/cosh^2 = sech^2; csch' = -cosh/sinh^2 = -csch coth; sech' = -sinh/cosh^2 = -sech tanh; coth' = (sinh^2-cosh^2)/sinh^2 = -csch^2
\begin{center}
\begin{tabular}[t]{|l|l|l|l|}
\hline
$f(x)$ & $f^{\prime}(x)$ & $f(x)$ & $f^{\prime}(x)$ \\
\hline
$\sinh (x)$ & $\answer{\cosh(x)}$ & $\csch(x)$ & $\answer{-\csch(x)\coth(x)}$ \\
\hline
$\cosh (x)$ & $\answer{\sinh(x)}$ & $\sech(x)$ & $\answer{-\sech(x)\tanh(x)}$ \\
\hline
$\tanh (x)$ & $\answer{\sech^{2}(x)}$ & $\coth(x)$ & $\answer{-\csch^{2}(x)}$ \\
\hline
\end{tabular}
\end{center}

Just like the trig functions, hyperbolic functions have inverses too.

\begin{center}
\begin{tabular}[t]{|l|l|l|}
\hline
$f(x)$ & $f^{-1}(x)$ & Domain \\
\hline
$\sinh (x)$ & $\arsinh(x)=\sinh ^{-1}(x)=\ln \left(x+\sqrt{x^{2}+1}\right)$ & $\mathbb{R}$ \\
\hline
$\cosh (x)$ & $\arcosh(x)=\cosh ^{-1}(x)=\ln \left(x+\sqrt{x^{2}-1}\right)$ & $[1, \infty)$ \\
\hline
$\tanh (x)$ & $\artanh(x)=\tanh ^{-1}(x)=\frac{1}{2} \ln \left(\frac{1+x}{1-x}\right)$ & $(-1, 1)$ \\
\hline
$\csch(x)$ & $\arcsch(x)=\csch^{-1}(x)=\ln \left(\frac{1}{x}+\sqrt{\frac{1}{x^{2}}+1}\right)$ & $\mathbb{R} \backslash\{0\}$ \\
\hline
$\sech(x)$ & $\arsech(x)=\sech^{-1}(x)=\ln \left(\frac{1+\sqrt{1-x^{2}}}{x}\right)$ & $(0,1]$ \\
\hline
$\coth(x)$ & $\arcoth(x)=\coth^{-1}(x)=\frac{1}{2} \ln \left(\frac{x+1}{x-1}\right)$ & $(-\infty,-1) \cup(1, \infty)$ \\
\hline
\end{tabular}
\end{center}

Similarly to the inverse trig functions, the inverse hyperbolic functions have derivatives.

\begin{center}
\begin{tabular}[t]{|l|l|l|l|}
\hline
$f(x)$ & $f^{\prime}(x)$ & $f(x)$ & $f^{\prime}(x)$ \\
\hline
$\arsinh(x)$ & $\frac{1}{\sqrt{1+x^{2}}}$ & $\arcsch(x)$ & $-\frac{1}{|x| \sqrt{x^{2}+1}}$ \\
\hline
$\arcosh(x)$ & $\frac{1}{\sqrt{x^{2}-1}}$ & $\arsech(x)$ & $-\frac{1}{x \sqrt{1-x^{2}}}$ \\
\hline
$\artanh(x)$ & $\frac{1}{1-x^{2}}$ & $\arcoth(x)$ & $\frac{1}{1-x^{2}}$ \\
\hline
\end{tabular}
\end{center}

\begin{warning}
$\artanh(x)$ and $\arcoth(x)$ appear to have the same function as a derivative, but recall that $\artanh(x)$ and $\arcoth(x)$ have different domains, so the derivatives are defined on different parts of the graph.
\end{warning}

\end{document}
